% ID: SEL-D0035-OA-E02
% ESTADO: validada. S036, cotejada con examen y criterios.
\textbf{a)} Las derivadas son $f'(x)=3x^2+2ax+b$ y
$f''(x)=6x+2a$. La inflexión en $-2$ exige $-12+2a=0$,
luego $a=6$. El mínimo en $-1$ exige $3-2a+b=0$,
de donde $b=9$. Con estos valores, $f''(-1)=6>0$ y
$f''$ cambia de signo en $-2$, por lo que se cumplen ambas
condiciones. Resulta $\boxed{a=6,\ b=9}$.

\textbf{b)} $f(x)=x(x+3)^2$, por lo que los cortes con el eje
horizontal son $(-3,0)$ y $(0,0)$. El corte vertical es $(0,0)$.
Como $f'(x)=3(x+3)(x+1)$, crece en $(-\infty,-3]$,
decrece en $[-3,-1]$ y vuelve a crecer en $[-1,\infty)$.
Hay un máximo relativo en $(-3,0)$ y un mínimo en $(-1,-4)$.
La raíz doble en $-3$ explica que la curva toque el eje sin
cruzarlo allí. La inflexión está en $(-2,-2)$.
\begin{center}
\setlength{\unitlength}{6mm}
\begin{picture}(10,9)(-5,-5)
\qbezier(-4.5,0)(-1.5,0)(1.5,0)\put(1.6,0){$x$}
\qbezier(0,-4.5)(0,-0.5)(0,3.5)\put(0,3.6){$y$}
\qbezier(-4,-4)(-3.75,-1.78125)(-3.5,-0.875)
\qbezier(-3.5,-0.875)(-3.25,0.03125)(-3,0)
\qbezier(-3,0)(-2.75,-0.03125)(-2.5,-0.625)
\qbezier(-2.5,-0.625)(-2.25,-1.21875)(-2,-2)
\qbezier(-2,-2)(-1.75,-2.78125)(-1.5,-3.375)
\qbezier(-1.5,-3.375)(-1.25,-3.96875)(-1,-4)
\qbezier(-1,-4)(-0.75,-4.03125)(-0.5,-3.125)
\qbezier(-0.5,-3.125)(-0.25,-2.21875)(0,0)
\qbezier(0,0)(0.125,1.12109)(0.25,2.64062)
\put(-3,0.3){\makebox(0,0)[b]{$(-3,0)$}}
\put(-1,-4.3){\makebox(0,0)[t]{$(-1,-4)$}}
\end{picture}
\end{center}
